% ============================================================
% RELATED WORK (critical evaluation; gaps; positioning)
% ============================================================
\section{Related Work}
\label{sec:related}

\subsection{Disparities in human pain care}
A consensus review assembled evidence that racial and ethnic minorities are
undertreated for pain across pain types (experimental, acute post-operative,
cancer, chronic non-malignant) and settings, including the emergency department,
and attributed the pattern to patient, provider and system factors
\cite{green2003unequal}; a later review traced its causes and consequences
\cite{anderson2009racial}, and a meta-analysis synthesised twenty years of
United States studies on analgesic treatment by race and ethnicity, examining
the moderating effect of pain type, setting and study period
\cite{meghani2012time}. Emergency-department opioid prescribing rose after the
quality initiatives of the late 1990s, but non-Hispanic White patients remained
more likely to receive an opioid than other groups \cite{pletcher2008trends}.
Hoffman et al.\ supplied a mechanism: false beliefs about biological differences
between Black and White people, endorsed by half of a sample of White medical
trainees, predicted lower pain ratings and less accurate treatment
recommendations for Black patients \cite{hoffman2016racial}. This literature
defines the disparity an LLM might inherit from clinical and lay text, namely
lower assessed severity and less potent analgesia for the same presentation, and
it fixes the direction that H2 predicts.

\subsection{Language-model audits in pain management}
Q-Pain introduced a dataset and an experimental framework for measuring social
bias in pain-management question answering and, applying it to GPT-2 and GPT-3,
found statistically significant treatment differences between intersectional
race--gender subgroups \cite{loge2021qpain}. Those models predate instruction
tuning and alignment, so the result establishes that the disparity is learnable
from text, not that current systems exhibit it. Young et al.\ brought the design
to GPT-4 and Gemini with 40 MIMIC-IV cases of abdominal, back, headache and
musculoskeletal pain, amended to every race/ethnicity and sex combination (480
cases per model), and found that race/ethnicity and sex did not influence
subjective pain rating, opioid name, order or dosage, although the two models
differed from each other in severity rating and strong-opioid use
\cite{young2025racial}; a subsequent letter proposed comparing inference-tuned
models with pretrained LLMs in this setting \cite{liu2025addressing}. In
palliative care, where pain management was one of four care dimensions probed,
GPT-4o produced expert-rated bias in a third of adversarial responses
\cite{akhras2026palliative}. Three features of the Young et al.\ study motivate
the present one: it tested proprietary models whose weights and updates are
outside the auditor's control; it stated race and sex explicitly; and its
univariate analyses did not pair responses on a sampling seed or report a
within-condition instability baseline, so the size of a difference that could
have been detected is unknown. The PAINT resource, with clinician-validated pain
and demographic annotations in four languages, names the question as open
\cite{zhang2026paint}.

\subsection{Demographic bias audits of clinical language models}
Beyond pain, the audits are numerous and largely concordant in finding bias,
but they differ in what they measure. Omiye et al.\ asked four commercial LLMs
nine questions probing race-based medical misconceptions, five times each, and
found that every model produced examples of debunked race-based medicine and
answered inconsistently across repetitions \cite{omiye2023race}. Zack et al.\
evaluated GPT-4 across medical education, diagnostic reasoning, clinical plan
generation and subjective patient assessment and found that it did not model the
demographic diversity of conditions and generated stereotyped vignettes
\cite{zack2024gpt4}. Hanna et al.\ fed 100 de-identified encounters of patients
with HIV to ChatGPT four times, varying only race/ethnicity, and compared the
generated text by sentiment, subjectivity and reading ease
\cite{hanna2023assessing}; Yang et al.\ found GPT-3.5 and GPT-4 projecting higher
costs and longer hospitalisations for White patients \cite{yang2024unmasking};
Ayoub et al.\ sampled simulated physicians choosing one patient to save under
resource scarcity \cite{ayoub2024inherent}. Benchmark-scale work followed:
Poulain et al.\ evaluated eight LLMs on three clinical question-answering
datasets with red-teaming prompts and examined how architecture and prompting
shape the biases observed \cite{poulain2024bias}, then proposed a
preference-optimisation alignment to reduce them \cite{poulain2024aligning};
FairMedQA built 4,806 counterfactual question pairs from 801 vignettes and found
accuracy disparities of 3 to 19 percentage points across demographic groups in
twelve LLMs \cite{xiao2025fairmedqa}; Purkayastha et al.\ found that
demographic prefixes changed 20.3\% of answers across six models on MedQA items
\cite{purkayastha2026beyond}; Lee et al.\ identified intersectional sex-by-race
gaps in LLM emergency triage \cite{lee2026promising}; and Pendse et al.\ altered
names, pronouns and descriptors across 39 demographic groups in 97 psychiatric
training cases to probe implicit bias in automated diagnosis
\cite{pendse2026implicit}. Audits of emergency triage that pair each case with
a counterfactual have found gender-swap flip rates between 9.9\% and 43.8\% in
five proprietary models \cite{young2026equitriage}, a dose-response skew toward
less aggressive management from stigmatising language
\cite{huang2026artificial}, and, when grounded in patient outcomes, harmful
reprioritisation from psychiatric history, frequent-visitor status, homelessness,
age and substance use but not from race, language or insurance
\cite{jarrett2026outcome}. Two systematic syntheses describe bias as pervasive
across clinical applications \cite{suenghataiphorn2025systematic,
xiao2026scoping}, and broader surveys catalogue the evaluation and mitigation
techniques available \cite{gallegos2024bias, mittermaier2023bias}.

Read critically, this body of work has three recurrent limitations that the
present study is designed around. Most audits test proprietary models through
an API, so results are tied to an unversioned endpoint. Nearly all state the
protected attribute explicitly, which is the cue most likely to trigger
alignment behaviour. And the common metric, the share of answers that change
when the attribute changes, is rarely compared with the share that change when
nothing changes.

\subsection{Open-weight models}
Open-weight models are now deployed on institutional hardware. Nowak et al.\
describe an isolation-first, on-premise inference stack for radiology that
served an open-weights model to 22 residents and radiologists in a prospective
pilot \cite{nowak2026secure}. Audits of such models are beginning: Aljohani
et al.\ compared ten open-source LLMs, including Qwen2.5, MedGemma, MedLLaMA2
and GPT-OSS variants, on paediatric Emergency Severity Index prediction with
paired counterfactuals that change one demographic, socioeconomic or
system-context variable \cite{aljohani2026counterfactual}; Ahmed and Sadka
audited Mistral 7B-Instruct, served locally through Ollama, on 4,659
counterfactual MedQA items under stochastic and deterministic decoding
\cite{ahmed2026auditing}; Kitahara and Yamaguchi found that job-posting language
triggers gender and racial bias in six open-weight models applied to recruitment,
among them Llama 3.2, Gemma 3 and Qwen 3 \cite{kitahara2026linguistic}; and
Bhetwal et al.\ measured consistent name bias in 300,000 responses from six local
LLMs on a fill-in-the-blank task \cite{bhetwal2026name}. The model families
audited here are documented in their technical reports \cite{gemmateam2025gemma,
yang2025qwen3, grattafiori2024llama}, and all are instruction-tuned in the
lineage of reinforcement learning from human feedback \cite{ouyang2022training}.
None has been audited for pain assessment.

\subsection{Name-based cues and audit instruments}
Name-based correspondence audits originate in labour economics, where otherwise
identical r\'esum\'es with White-sounding names received more call-backs than
those with Black-sounding names \cite{bertrand2004emily}. Applied to LLMs, the
design shows that advice systematically disadvantages names associated with
racial minorities and women, with Black women receiving the least advantageous
outcomes across 42 prompt templates \cite{salinas2024name}; that hiring emails
favour White over Hispanic applicants, although the pattern is prompt-sensitive
and idiosyncratic \cite{an2024hiring}; that hiring and salary recommendations
over 750,000 prompts favour White female-sounding names \cite{nghiem2024gotta};
and that a method varying demographics across 70 decision scenarios can surface
discrimination before deployment \cite{tamkin2023evaluating}. Cues beyond the
name matter too: models make covertly racist decisions from dialect alone
\cite{hofmann2024dialect}, and demographic identity leaks from unstructured
prose in de-identified r\'esum\'es across nine open-weight models
\cite{chen2026beyond}. Two recent results bear directly on how such audits
should be built. Tonneau et al.\ showed over 14.8 million prompts that different
cues for the same group induce only partially overlapping changes in model
responses, so that bias conclusions vary in magnitude and direction with the
cue chosen \cite{tonneau2026cues}. Vohra and Ravikiran, in a preregistered
replication across hiring, lending and medical triage with applications that
differed only by name, found that none of 36 planned contrasts survived
correction and that the audit instrument (rating versus ranking) produced
effects rivalling the demographic ones \cite{vohra2026audit}.

\subsection{Instability as a baseline}
Atil et al.\ documented accuracy variation of up to 15\% across runs of LLMs
configured to be deterministic \cite{atil2024nondeterminism}, and Song et al.\
argued that evaluation should not ignore non-determinism \cite{song2025good}.
For fairness audits specifically, Bellibatlu et al.\ measured how often a
clinical LLM agent changes its action when the prompt is re-run unchanged: 8.7\%
of replicate pairs at default sampling for one agent, from 2.2\% for
intensive-care escalation to 17.9\% for controlled-substance caution, and pooled
floors of 2.5\% to 23.7\% across six models from five vendors
\cite{bellibatlu2026instability}. A counterfactual flip rate below that floor
cannot be attributed to the demographic change, and controlled-substance
decisions sat at the unstable end of their range, which is directly relevant to
an opioid audit.

\subsection{Position of the present study}
To the authors' knowledge this is the first audit of open-weight, locally
deployable models for racial and sex disparities in pain assessment; the first
in this task to cue identity by name alone with two names per group, so that a
name's idiosyncrasy can be separated from a group effect; and the first to pair
every counterfactual comparison on the sampling seed and to report the
re-sampling noise floor and the minimum detectable effect alongside the
estimates. It extends the explicit-label design of Young et al.\
\cite{young2025racial} to open-weight models, imports the name-based audit
design from labour economics and general LLM auditing
\cite{bertrand2004emily, salinas2024name} into pain assessment, and adopts the
instability baseline of Bellibatlu et al.\ \cite{bellibatlu2026instability}.
